% ============================================================
% 070_T0_Mathematische_Struktur_De_ch.tex
% Zur mathematischen Struktur der T0-Theorie
% Warum Produkte verschiedener r_i strukturell bedeutungslos sind
% Johann Pascher — Mai 2026
% ============================================================

\section*{Vorbemerkung}

Die im Folgenden diskutierten Operationen sind \textbf{algebraisch korrekt}.
Die Einschränkungen sind keine Rechenregeln, sondern
\textbf{Konsistenzbedingungen des FFGFT-Modells}:
Produkte verschiedener $r_i$ sind algebraisch definiert,
aber strukturell bedeutungslos --- sie entsprechen keinem
Zustand im Modell.

%------------------------------------------------------------
\section{Die Vorfaktoren als geometrische Invarianten}
%------------------------------------------------------------

In FFGFT gilt für jedes Teilchen $i$:
\begin{equation}
  m_i = r_i \cdot \xi_0^{p_i} \cdot v,
  \qquad r_i = f(\mathcal{T}_i), \quad p_i = g(\mathcal{T}_i)
\end{equation}
wobei $\mathcal{T}_i = (n_i, l_i, j_i)$ die Torusquantenzahl des Teilchens ist.

Die Funktionen $f$ und $g$ sind \textbf{injektiv} auf der Menge der
zulässigen Torusmoden: Aus $(r_i, p_i)$ kann man $\mathcal{T}_i$ eindeutig
rekonstruieren.

Die Vorfaktoren $r_i$ sind \textbf{keine freien Parameter}, sondern
\textbf{geometrische Invarianten} des jeweiligen Teilchens ---
vollständig durch dessen Quantenzahlen $(n,l,j)$ festgelegt.

Für Elektron und Myon (Dok.~006):
\begin{align}
  m_e &= \frac{4}{3} \cdot \xi_0^{3/2} \cdot v \label{eq:me}\\
  m_\mu &= \frac{16}{5} \cdot \xi_0^{1} \cdot v \label{eq:mmu}
\end{align}

%------------------------------------------------------------
\section{Die Menge der Paare \texorpdfstring{\texorpdfstring{\texorpdfstring{$(r_i,p_i)$}{(ri,pi)}}{(ri,pi)}}{(ri,pi)} ist multiplikativ nicht abgeschlossen}
%------------------------------------------------------------

\subsection{Die Torus-Quantisierung}

Die erlaubten Torusmoden erzeugen $r_i$ und $p_i$ durch spezifische
geometrische Regeln (Dok.~006). Die zulässigen Paare $(r,p)$
sind diskret und endlich:

\begin{center}
{\footnotesize
\begin{tabular}{lcccc}
\toprule
\textbf{Teilchen} & \textbf{$(n,l,j)$} & \textbf{$r_i$} & \textbf{$p_i$} \\
\midrule
Elektron & $(1,0,\nicefrac{1}{2})$ & $4/3$ & $3/2$ \\
Myon     & $(2,1,\nicefrac{1}{2})$ & $16/5$ & $1$ \\
Tau      & $(3,2,\nicefrac{1}{2})$ & $25/9$ & $2/3$ \\
\bottomrule
\end{tabular}
}
\end{center}

\subsection{Das Nichtabschluss-Theorem}

Bildet man das Produkt $r_e \cdot r_\mu$:
\[
  r_e \cdot r_\mu = \frac{4}{3} \cdot \frac{16}{5} = \frac{64}{15},
  \qquad p_e + p_\mu = \frac{3}{2} + 1 = \frac{5}{2}
\]

Das Modell fragt:
\begin{center}
  \emph{Gibt es ein Teilchen $X$ mit $r_X = 64/15$ und $p_X = 5/2$?}
\end{center}

\textbf{Antwort: Nein.} In der Quantenzahltabelle kommt $64/15$ bei
keinem $p$ vor, und $p=5/2$ bei keinem $r$.

Die Menge der zulässigen Paare $\mathcal{P}$ ist bezüglich der
Multiplikation \textbf{nicht abgeschlossen}:
\begin{equation}
  \forall\, i \neq j:\quad (r_i \cdot r_j,\; p_i + p_j) \notin \mathcal{P}
  \label{eq:nichtabschluss}
\end{equation}

\subsection{Warum ist das so?}

Die erlaubten Torusmoden haben separate, nicht multiplikative
Erzeugungsregeln für $r$:
\[
  r(n,l) = \frac{(l+1)^2}{n^2 - l(l+1)} \cdot \text{(Spin-Anteil)}
\]
(vereinfacht). Das Produkt zweier solcher Brüche ist \textbf{nie}
wieder ein Bruch derselben Form mit zwei Quantenzahlen.

Die $r_i$ sind daher keine Zahlen in einem Körper, sondern
\textbf{Etiketten eines diskreten geometrischen Raums mit eigener
Verknüpfung} --- und diese Verknüpfung ist nicht die Körpermultiplikation.

%------------------------------------------------------------
\section{Algebraisch korrekt, aber strukturell bedeutungslos}
%------------------------------------------------------------

\subsection{Das Kreuzprodukt-Problem}

Setzt man \eqref{eq:me} und \eqref{eq:mmu} symbolisch in
$E_0 = \sqrt{m_e \cdot m_\mu}$ ein:
\begin{equation}
  E_0^2 = m_e \cdot m_\mu
        = \underbrace{\frac{4}{3} \cdot \frac{16}{5}}_{r_e \cdot r_\mu}
          \cdot \xi_0^{3/2 + 1} \cdot v^2
        = \frac{64}{15} \cdot \xi_0^{5/2} \cdot v^2
  \label{eq:kreuz}
\end{equation}

Das Ergebnis ist \textbf{algebraisch korrekt}. Aber weil
$(64/15,\;5/2) \notin \mathcal{P}$, entspricht dieser Ausdruck
\textbf{keinem FFGFT-Teilchen} --- er ist strukturell nicht
interpretierbar.

\subsection{Analogie: verschiedene algebraische Räume}

Jedes Teilchen $i$ lebt in einem durch $\mathcal{T}_i$ aufgespannten
Zustandsraum. Die $r_i$ sind Koeffizienten in \emph{verschiedenen Räumen}.
Ihr Produkt entspricht einem Tensorprodukt:
\[
  r_e \cdot r_\mu \;\leftrightarrow\; \mathbf{v}_e \otimes \mathbf{v}_\mu
\]
Das ergibt einen Tensorprodukt-Zustand --- \textbf{nicht} einen Zustand
im ursprünglichen Ein-Teilchen-Raum. In einer Ein-Teilchen-Theorie
(Massenformeln) ist das bedeutungslos.

\subsection{Weitereinsetzen verschlechtert das Ergebnis}

Setzt man \eqref{eq:kreuz} in $\alpha = \xi_0 \cdot (E_0/\text{MeV})^2$:
\begin{equation}
  \alpha \stackrel{\text{strukturell falsch}}{=}
  \frac{64}{15} \cdot \xi_0^{7/2} \cdot \frac{v^2}{\text{MeV}^2}
\end{equation}
Exponent $7/2$ und Faktor $64/15$ kommen nirgendwo in den
Grundgleichungen vor. Zusätzlich fehlt $K_{\text{frak}}$,
und Einheiten werden verdeckt vermischt.
In der Praxis: statt $-1{,}35\,\%$ ergibt sich $-3{,}15\,\%$
Abweichung.

%------------------------------------------------------------
\section{Was erlaubt ist und warum}
%------------------------------------------------------------

\subsection{Numerische Auswertung: Übergang nach $\mathbb{R}$}

Zahlen leben in $\mathbb{R}$, wo Multiplikation immer erlaubt ist.
Sobald $r_e$ und $r_\mu$ zu Zahlen ausgewertet wurden, haben sie
\textbf{keine Erinnerung} mehr an ihre Herkunft:

\begin{align*}
  m_e &\approx 0{,}511\,\text{MeV}, \quad
  m_\mu \approx 105{,}66\,\text{MeV} \\
  E_0 &= \sqrt{0{,}511 \cdot 105{,}66}\,\text{MeV} = 7{,}348\,\text{MeV} \\
  \alpha &= \xi_0 \cdot (7{,}348)^2 \approx 0{,}00720 \quad (-1{,}35\,\%)
\end{align*}

\subsection{Division: Respektieren der getrennten Räume}

Bei Verhältnissen kürzen sich $v$ und $K_{\text{frak}}$ heraus:
\[
  \frac{m_\mu}{m_e}
  = \frac{r_\mu}{r_e} \cdot \xi_0^{p_\mu - p_e}
  = \frac{12}{5} \cdot \xi_0^{-1/2}
  \approx 207{,}9 \quad (0{,}5\,\%)
\]
Das Ergebnis $12/5$ ist interpretierbar --- kein Kreuzprodukt.
Die Division \textbf{respektiert} die getrennten Räume.

\begin{tcolorbox}[colback=blue!5, colframe=blue!60!black, boxrule=0.5pt,
  arc=2pt, left=3mm, right=3mm, fontupper=\small,
  title={Grundregel --- Die Nichtabschluss-Bedingung}]
Die Menge der zulässigen Paare $(r_i, p_i)$ ist bezüglich der
Multiplikation nicht abgeschlossen:
\[
  \boxed{\forall\, i \neq j:\; (r_i \cdot r_j,\; p_i + p_j) \notin \mathcal{P}}
\]
Daher: Produkte verschiedener $r_i$ sind algebraisch definiert,
entsprechen aber keinem Teilchen im Modell.
Erst numerische Auswertung (Übergang nach $\mathbb{R}$) erlaubt
freie Kombination.
\end{tcolorbox}

%------------------------------------------------------------
\section{Verbindung zu anderen Resultaten}
%------------------------------------------------------------

Dasselbe Prinzip gilt für die Koide-Relation (Dok.~192):
Auch dort entstehen beim symbolischen Einsetzen Kreuzterme mit
Exponenten, die zu keinem Teilchen gehören.

Die Regel ist keine Willkür, sondern eine \textbf{notwendige
Konsequenz der Torus-Quantisierung} und der Injektivität von
$f$ und $g$.

%------------------------------------------------------------
\section*{Zusammenfassung}
%------------------------------------------------------------

\begin{quote}
\emph{Die Vorfaktoren $r_i$ sind keine freien Parameter, sondern
geometrische Invarianten von Torusmoden mit festen Quantenzahlen.
Die Menge der zulässigen Paare $(r_i, p_i)$ ist bezüglich der
Multiplikation nicht abgeschlossen --- ein Produkt $r_i \cdot r_j$
mit $i \neq j$ ist algebraisch definiert, aber keiner Torusmode zugehörig
und damit im Modell bedeutungslos. Kombinationen sind erst nach
numerischer Auswertung in $\mathbb{R}$ erlaubt.}
\end{quote}

%------------------------------------------------------------
\section*{Zeichenerklärung}
%------------------------------------------------------------

\begin{center}
{\footnotesize
\begin{tabular}{lp{7cm}}
\toprule
\textbf{Symbol} & \textbf{Bedeutung} \\
\midrule
$\mathcal{T}_i = (n_i, l_i, j_i)$ & Torusmode: Haupt-, Bahn-, Spinquantenzahl \\
$r_i = f(\mathcal{T}_i)$ & Geometrischer Vorfaktor (injektiv) \\
$p_i = g(\mathcal{T}_i)$ & Exponent in der Massenformel \\
$\mathcal{P}$ & Menge der zulässigen Paare $(r_i, p_i)$ \\
$v = 246\,\text{GeV}$ & Higgs-Vakuumerwartungswert \\
$\xi_0 = 4/30000$ & Geometrischer Basisparameter \\
$K_{\text{frak}}$ & Fraktale Korrektur ($\approx 0{,}986$) \\
$\mathbb{R}$ & Reelle Zahlen (keine Strukturerinnerung) \\
$\mathbf{v}_e \otimes \mathbf{v}_\mu$ & Tensorprodukt (in Ein-Teilchen-Theorie bedeutungslos) \\
$(64/15,\;5/2)$ & Kreuzpaar $\notin \mathcal{P}$, kein FFGFT-Teilchen \\
\bottomrule
\end{tabular}
}
\end{center}


\medskip\noindent
\textbf{Verweise:} Dok.~006 (Quantenzahlen), Dok.~043, Dok.~011,
Dok.~044 ($\alpha$), Dok.~192 (Koide, allgemeines Prinzip).
